ابتدا با توجه به اطلاعات داده شده در بخش‌های قبل، ویژگی‌های پاسخ ضربه سیستم یعنی
$h[n]$
را مشخص می‌کنیم. از شرط اول و دوم می‌دانیم سیستم علی است
($h[n]=0$ برای $n<0$)
و
$h[n]$
دنباله‌ای با مقادیر حقیقی است. از شرط سوم که بیان می‌کند
\lr{DTFT}
سیگنال
$h[n+1]$
حقیقی است، نتیجه می‌گیریم که این دنباله باید متقارن زوج باشد:
\begin{flalign*}
    h[n+1] = h[-n+1] \implies h[m] = h[2-m] &&
\end{flalign*}
با در نظر گرفتن شرط علیت، دنباله
$h[n]$
تنها برای
$n=0, 1, 2$
می‌تواند غیر صفر باشد و از رابطه تقارن داریم
$h[0] = h[2]$.

حال برای حل بخش چهارم، از قضیه پارسوال
(\lr{Parseval's theorem})
استفاده می‌کنیم. این قضیه بیان می‌کند که انرژی سیگنال در حوزه زمان با انتگرال انرژی در حوزه فرکانس برابر است:
\begin{flalign*}
    \sum_{n=-\infty}^{\infty} |h[n]|^2 = \frac{1}{2\pi} \int_{-\pi}^{\pi} |H(e^{j\omega})|^2 d\omega = 2 &&
\end{flalign*}
با جایگذاری مقادیر غیر صفر
$h[n]$
در بسط رابطه بالا خواهیم داشت:
\begin{flalign*}
    h[0]^2 + h[1]^2 + h[2]^2 = 2 &&
\end{flalign*}
از آنجا که پیش‌تر نشان دادیم
$h[0] = h[2]$
است، می‌توانیم رابطه را به فرم ساده‌تر زیر بنویسیم:
\begin{flalign*}
    2h[0]^2 + h[1]^2 = 2 &&
\end{flalign*}
بنابراین پاسخ ضربه سیستم به فرم
$h[n] = a\delta[n] + b\delta[n-1] + a\delta[n-2]$
است که ضرایب آن
($a=h[0]$ و $b=h[1]$)
در معادله به دست آمده صدق می‌کنند.